\section*{Tóm tắt nội dung và đóng góp}

Báo cáo trình bày Knowledge Hub, một dịch vụ GraphRAG đóng vai trò bộ nhớ dài hạn và nguồn cung cấp ngữ cảnh cho các AI Agent trong quy trình phát triển phần mềm của một team. Hệ thống thu thập tri thức hỗn hợp của một product -- mã nguồn, tài liệu, lịch sử git và đặc tả -- cấu trúc hóa bằng cách kết hợp vector embedding với knowledge graph, và phục vụ truy xuất hybrid (ngữ nghĩa kết hợp từ khóa và duyệt đồ thị) qua hai giao diện MCP và REST, với tri thức dùng chung cho cả team và kiểm soát truy cập ở mức nguồn.

Về nội dung, báo cáo trình bày trọn vòng đời của tri thức trong hệ thống: cơ chế định danh ổn định làm nền tảng cho đồng bộ incremental; chiến lược chunking giữ trọn ngữ nghĩa theo từng loại dữ liệu; việc dựng và liên kết đồ thị tri thức có phân biệt liên kết chắc chắn với liên kết phỏng đoán qua điểm tin cậy; truy xuất hybrid hợp nhất nhiều đường tìm kiếm theo thứ hạng; và phân quyền ACL áp như một bộ lọc cứng trên mọi đường truy xuất. Toàn hệ thống được đóng gói bằng Docker và cấu hình được theo nhiều môi trường.

Về kết quả, hệ thống đã hiện thực đầy đủ các tính năng trong phạm vi và đạt các tiêu chí phi chức năng đề ra ở quy mô đánh giá: độ trễ truy vấn p95 khoảng $254$ ms khi chưa cache và khoảng $5$ ms khi đã cache; cập nhật tăng dần hoàn tất dưới hai giây cho một thay đổi nhỏ; chi phí độ trễ của phân quyền không đáng kể (khoảng $1\%$); và chất lượng truy xuất đạt Recall@10 khoảng $0.96$ với MRR khoảng $0.75$ trên bộ truy vấn có nhãn, trong đó truy vấn hybrid cải thiện thứ hạng kết quả so với tìm kiếm ngữ nghĩa thuần. Bộ kiểm thử tự động chạy trên mỗi thay đổi bảo đảm tính đúng đắn về chức năng và không rò rỉ phân quyền.

Đóng góp chính của đề tài, bên cạnh một hệ thống hoàn chỉnh chạy được, là làm rõ ranh giới phân vai giữa hệ thống và agent -- hệ thống \textit{truy xuất} và cung cấp ngữ cảnh kèm căn cứ, còn agent \textit{suy luận} -- cùng một kiến trúc phân tầng theo cổng và bộ chuyển đổi cho phép mở rộng (thêm ngôn ngữ, thêm loại quan hệ, thêm nguồn) bằng cách bổ sung thành phần thay vì sửa phần lõi.
